% from main.tex
%%%%%%%%%%  scan idx 95  |  folio 82  %%%%%%%%%%
sont donnés dans l'Exp.\ II.

On montre en particulier qu'il existe des résolutions globales dans les cas suivants : \(X\) est le topos zariskien d'un schéma affine, ou plus généralement d'un schéma quasi-compact possédant un faisceau inversible ample, voire une famille ample de faisceaux inversibles au sens de Kleiman (par exemple un schéma régulier); ou encore \(X\) est le topos des faisceaux d'ensembles sur un espace topologique compact, \(\mathcal O_X\) le faisceau des fonctions continues à valeurs complexes. À titre d'illustration~-- et pour mettre en évidence la souplesse des notions introduites~-- on donne en appendice une définition «~purement faisceautique~» de l'indice d'une famille d'opérateurs elliptiques.

L'Exp.\ III étudie la stabilité des conditions de finitude par image directe dérivée. Pour obtenir des énoncés raisonnables, il y a lieu de «~relativiser~» les notions de pseudo-cohérence et de perfection. On se place ici dans le cadre des schémas ordinaires, ce qui suffit pour le Séminaire, mais il ne fait pas de doute qu'il faudra tôt ou tard développer une théorie analogue pour les schémas ou espaces analytiques relatifs. Soit \(p:X\to S\) un \(S\)-schéma localement de type fini et soit \(F\) un complexe de \(\mathcal O_X\)-Modules, on dit que \(F\) est \emph{pseudo-cohérent} (\emph{resp.\ parfait}) \emph{relativement à} \(p\) si l'on peut localement plonger (par une immersion fermée) \(X\) dans un \(S\)-schéma lisse \(X'\) de manière que \(F\) prolongé par zéro sur \(X'\) soit pseudo-cohérent (resp.\ parfait). La notion de pseudo-cohérence relativement à \(p\) est surtout intéressante dans le cas où \(S\) n'est pas localement noethérien, car dans le cas contraire elle coïncide avec la notion de pseudo-cohérence ordinaire. D'autre part, comme on peut s'y attendre, dire que \(F\) est parfait rel.\ à \(p\)
